% Source: physical PDF pp. 178--180 (printed pp. 155--157), from the Section
% 27.7 heading on physical p. 178 through "may be considered to be soft." on
% physical p. 180, immediately before the Section 27.8 heading.
% Inventory: Eqs. (27.7.1)--(27.7.4), no unnumbered displays, four linked
% citation occurrences referring to References 8--10, one ordinary
% section-title footnote, and no figures, tables, or typographic dividers.
% Uncertainties: none.

\subsection[Soft Supersymmetry Breaking]{Soft Supersymmetry Breaking\footnotemark}
\footnotetext{This section lies somewhat out of the book's main line of
development and may be omitted in a first reading.}

We will see in the next chapter that, even if supersymmetry is an exact
symmetry of the action, the spontaneous breakdown of supersymmetry at very
high energy can produce superrenormalizable terms that violate supersymmetry
conservation in the effective action that describes physics at lower
energies. These superrenormalizable terms may explain the lack of
supersymmetry observed in phenomena at accessible energies. In this section
we will consider the radiative corrections that can be produced by such
supersymmetry-breaking superrenormalizable terms, in part to see whether this
provides a criterion for including or rejecting such terms in supersymmetric
versions of the standard model.

The sign of supersymmetry breaking is the appearance of expectation values
of \(D\)-terms of general superfields or \(\mathcal{F}\)-terms of chiral
superfields. Any operator \(\epsilon\mathcal{O}\) in the Lagrangian density
that breaks supersymmetry can be written in a supersymmetric form, as a
\(D\)-term
\begin{equation}
  \epsilon\mathcal{O}=[ZS]_D\,,
  \tag{27.7.1}\label{eq:27.7.1}
\end{equation}
where \(S\) is a non-chiral superfield that has \(\mathcal{O}\) as its
\(C\)-term, and \(Z\) is a non-chiral external superfield whose only
non-vanishing component is \([Z]_D=\epsilon\). Some but not all operators
\(\epsilon\mathcal{O}\) that break supersymmetry can also be written as
\(\mathcal{F}\)-terms,
\begin{equation}
  \epsilon\mathcal{O}=[\Omega O]_{\mathcal{F}}\,,
  \tag{27.7.2}\label{eq:27.7.2}
\end{equation}
or their adjoints, where \(O\) is the left-chiral superfield whose
\(\mathcal{F}\)-term is \(\mathcal{O}\) and \(\Omega\) is an external
left-chiral superfield whose only non-vanishing component is
\([\Omega]_{\mathcal{F}}=\epsilon\). We can count the order in \(\epsilon\)
in which a given correction to the effective Lagrangian will occur by
counting the powers of \(Z\) or \(\Omega\) needed to construct this correction
in a supersymmetric way. We will find interesting limitations on the
radiative corrections that can be produced by those interactions that can be
written in the form \eqref{eq:27.7.2} as well as \eqref{eq:27.7.1}.

According to the results of the previous section, there are no radiative
corrections to \(\mathcal{F}\)-terms, so all supersymmetry-breaking radiative
corrections to the Wilsonian Lagrangian density must take the form of
\(D\)-terms. This theorem does not prevent any given operator from appearing
in the Wilsonian Lagrangian density, because even if an operator
\(\epsilon\Delta\mathcal{L}\) cannot be expressed in the form
\([Z\Lambda]_D\), where \(\Lambda\) is the general superfield whose
\(C\)-term is \(\Delta\mathcal{L}\), yet \(\epsilon^2\Delta\mathcal{L}\)
can be expressed in the form
\begin{equation}
  \epsilon^2\Delta\mathcal{L}
  =2[\Omega^*\Omega\Lambda]_D\,.
  \tag{27.7.3}\label{eq:27.7.3}
\end{equation}
But not all operators can be produced by radiative corrections of
\emph{first order} in \(\Omega\) or \(\Omega^*\). In particular, a function
only of the \(\phi\)-term of a left-chiral superfield \(\Phi\), but not of
\(\phi^*\), cannot be written as the \(D\)-term of a superfield linear in
\(\Omega\). (Note that \([\Omega h(\Phi)]_D\) is a derivative, while
\([\Omega^*h(\Phi)]_D
=2[\Phi]_{\mathcal{F}}\,\partial h(\phi)/\partial\phi\) is not a function
of \(\phi\) alone.) We conclude then that \emph{the
supersymmetry-breaking terms in the Wilsonian Lagrangian that depend on
\(\phi\) alone cannot be produced by radiative corrections that are of first
order in supersymmetry-breaking interactions of the form
\eqref{eq:27.7.2}.}

This result is significant because the most divergent radiative corrections
are those that are of lowest order in superrenormalizable couplings. To be
more specific, the coefficient of an interaction of dimensionality
\(\mathcal{D}\) has dimensionality (in powers of energy)
\(4-\mathcal{D}\), so dimensional analysis indicates that the contribution
of a set of interactions of dimensionality \(d_1\), \(d_2\), etc. to the
coefficient of an interaction of dimensionality \(d\) can contain the
ultraviolet cut-off to at most a power
\begin{equation}
  p=4-d-(4-d_1)-(4-d_2)-\cdots
  \tag{27.7.4}\label{eq:27.7.4}
\end{equation}
and is therefore finite if \(p<0\). (This argument ignores possible
ultraviolet divergences in subintegrations; for a thorough treatment of this
topic, see Reference~\hyperref[ch27-ref-8]{8}.) Superrenormalizable
interactions are `soft,' in the sense that they reduce the degree of
divergence of the graphs in which they appear. In particular, in a
renormalizable theory where all interactions have \(d_i\leq4\), and the
strictly renormalizable interactions with \(d_i=4\) are supersymmetric, the
contribution of one or more superrenormalizable interactions to the
coefficient of an interaction with \(d=4\) will always have \(p<0\), so even
if they are not supersymmetric the superrenormalizable interactions will not
produce supersymmetry-violating ultraviolet-divergent corrections to the
coefficients of the supersymmetric \(d=4\) interactions.

On the other hand, in such a theory there may be divergent radiative
corrections to the superrenormalizable interactions
themselves~\hyperref[ch27-ref-9]{[9]}. The most worrisome are quadratic (or
higher) divergences, which if cut off at some high energy scale \(M_X\) may
require a fine-tuning of the bare coupling constants to preserve
supersymmetry as a good approximate symmetry at energies below \(M_X\).
According to Eq.~\eqref{eq:27.7.4}, in renormalizable theories where all
interactions with \(d_i=4\) are supersymmetric, radiative corrections can
produce quadratic or more highly divergent (\(p\geq2\))
supersymmetry-violating operators of dimensionality \(d\) only if they
involve insertion of a superrenormalizable supersymmetry-violating
interaction of dimensionality \(d_1\geq2+d\). This allows either \(d=0\)
and \(d_1\geq2\), which arises only when we calculate the cosmological
constant, or \(d=1\) and \(d_1=3\), which arises only when we calculate the
`tadpole' graphs in which a scalar field line disappears into the vacuum.
The cosmological constant raises fine-tuning problems for all known
theories~\hyperref[ch27-ref-10]{[10]} and will not be considered further
here. The tadpole graphs represent operators linear in \(\phi\) or \(\phi^*\)
and, as we have seen, can not be produced to first order in
supersymmetry-violating interactions that can be put in the form
\eqref{eq:27.7.2}. Such superrenormalizable interactions are therefore
`soft' in the sense that they do not induce quadratic or higher divergences.
Along with the superrenormalizable interactions with \(d\leq2\), including
arbitrary quadratic polynomials in \(\phi\) and \(\phi^*\), the
supersymmetry-breaking interactions that are soft in this sense include
terms of third order in the \(\phi\), which can be expressed as
\(\phi^3=[\Omega\Phi^3]_{\mathcal{F}}\), and likewise terms of third order
in the \(\phi^*\), and also \(d=3\) gaugino mass terms, which can be written
as \([\Omega\epsilon_{\alpha\beta}W_\alpha W_\beta]_{\mathcal{F}}\), but
not terms like \(\phi^2\phi^*\) or \(\phi\phi^{2*}\), which in general
\emph{can} produce quadratically divergent tadpole
graphs~\hyperref[ch27-ref-9]{[9]}.

Nevertheless, tadpoles can only arise for scalar fields that are neutral
with respect to all exact symmetries. In theories without such neutral
scalars, like the supersymmetric standard model discussed in the next
chapter, \emph{all} superrenormalizable interactions may be considered to be
soft.
